\documentclass[12pt,a4paper]{article}

% --- Paquetes de diseño y fuentes ---
\usepackage[utf8]{inputenc}
\usepackage[spanish, es-tabla]{babel}
\usepackage{amsmath, amsfonts, amssymb, amsthm}
\usepackage{graphicx}
\usepackage{geometry}
\usepackage{hyperref}
\usepackage{booktabs}
\usepackage{fancyhdr}
\usepackage{setspace}

% Configuración de márgenes
\geometry{top=2.5cm, bottom=2.5cm, left=2.5cm, right=2.5cm}
\onehalfspacing

% Cabecera y Pie de página
\pagestyle{fancy}
\fancyhf{}
\rhead{\small Semillero de Aprendizaje por Refuerzo}
\lhead{\small IA para Trading y Sentimiento}
\cfoot{\thepage}
\renewcommand{\headrulewidth}{0.4pt}

% Título y Autores
\title{
    \vspace{-1cm}
    \textbf{\Large IA y Emociones: Desarrollo de un Agente de Aprendizaje por Refuerzo para el Mercado de Criptomonedas} \\
    \vspace{0.5cm}
    \large Universidad Externado de Colombia \\
    \text{Ciencia de Datos -- Departamento de Matemáticas}
}
\author{
    \textbf{Julian Duarte} \quad \textbf{Julian Jimenez} \\
    \medskip
    \small \textit{Tutores:} \\
    \small Prof. Julián Fernando Sánchez \quad Prof. Juan Carlos Urueña \quad Prof. Camilo Castillo
}
\date{Marzo de 2026}

\begin{document}

\maketitle

\begin{abstract}
    Este trabajo presenta la creación de una Inteligencia Artificial (Agente de Aprendizaje por Refuerzo) capaz de operar automáticamente en el mercado de criptomonedas. A diferencia de los modelos tradicionales que solo miran números, nuestra propuesta integra el "estado de ánimo" del mercado (sentimiento colectivo) para tomar mejores decisiones. Este documento explica paso a paso el diseño matemático y lógico detrás de este agente, definiendo cómo percibe el mundo, qué acciones puede tomar y qué beneficios espera obtener.
\end{abstract}

\section{Problema: Mercado y Datos}
El desarrollo de este proyecto se centra en optimizar la toma de decisiones en el trading (compra y venta de activos), una tarea compleja debido a la alta volatilidad y el componente emocional de los mercados. Operaremos en el intercambio \textbf{Binance} usando el par \textbf{BNB/USDT} (comprar Binance Coin usando Tether, una "stablecoin" o moneda estable ligada al valor del dólar). 

Para que nuestra IA aprenda, utilizamos dos tipos de información fundamental:
\begin{itemize}
    \item \textbf{Datos OHLCV (Precios):} Corresponde a los precios de mercado (Apertura, Máximo, Mínimo, Cierre) y Volumen. Los usamos para calcular el \textbf{retorno logarítmico}: $r_t = \log(P_t/P_{t-1})$.
    \item \textbf{Sentimiento de Mercado:} Integramos índices que miden si los inversores tienen "Miedo" o "Codicia". Esta variable psicológica es vital en el ecosistema cripto.
\end{itemize}

\subsection{Alternativas para la Fuente de Sentimiento}
Para medir este componente emocional, planteamos tres posibles fuentes tecnológicas:
\begin{enumerate}
    \item \textbf{Redes Sociales Crudas (X y Reddit):} Análisis directo de publicaciones masivas mediante técnicas de Procesamiento de Lenguaje Natural (NLP).
    \item \textbf{Comunidades Especializadas (Stocktwits):} Plataforma donde los usuarios etiquetan sus mensajes como alcistas o bajistas de forma explícita.
    \item \textbf{Índices Agregadores (Crypto Fear \& Greed Index):} Valor diario consolidado basado en múltiples indicadores de mercado y redes sociales.
\end{enumerate}

\section{Espacio de Estados: ¿Qué observa la IA?}
El "Estado" es la información que la IA procesa antes de decidir qué acción tomar. Definimos el estado $S_t = (M_t, I_t, Sent_t)$:
\begin{enumerate}
    \item \textbf{Régimen de Mercado ($M_t$):} Clasificamos el mercado como alcista (\textbf{U}p) o bajista (\textbf{D}own) mediante $r_t > \theta$.
    \item \textbf{Inventario ($I_t$):} Indica la posición actual del agente ($0$: sin posición, $1$: largo o con el activo comprado). Operamos en el mercado \textbf{Spot} (intercambio real e inmediato de activos).
    \item \textbf{Sentimiento de Mercado ($Sent_t$):} Nivel de optimismo o pesimismo colectivo.
\end{enumerate}

\section{Espacio de Acciones y Admisibilidad}
El agente elige entre $\mathcal{A} = \{Buy, Sell, Hold\}$. La admisibilidad depende de $I_t$:
\begin{itemize}
    \item Si $I_t=0$ (sin activo), puede elegir entre $\{Buy, Hold\}$.
    \item Si $I_t=1$ (con activo), puede elegir entre $\{Sell, Hold\}$.
\end{itemize}

\section{Probabilidades de Transición y Estimación}
Para modelar cómo cambia el mercado, estimamos las probabilidades de transición mediante el **Suavizado de Laplace**, evitando nulidad en eventos no observados:
\begin{equation}
    \tilde{\alpha} = \frac{N_{UU} + \kappa}{N_{U \bullet} + 2\kappa}, \quad \tilde{\beta} = \frac{N_{DD} + \kappa}{N_{D \bullet} + 2\kappa}
\end{equation}

\section{Recompensas: Beneficios y Costos}
La recompensa $R_{t+1}$ se calcula como la ganancia bruta del precio menos los **costos operativos** ($c$), que incluyen comisiones y el \textbf{Spread} (diferencia de precios de ejecución):
\begin{equation}
    R_{t+1} = I_{t+1}r_{t+1} - c|I_{t+1} - I_t|
\end{equation}
La recompensa esperada pondera los escenarios futuros del mercado mediante la media condicional $\bar{\mu}(m)$:
\begin{equation}
    \bar{\mu}(U) = \alpha \mu_U + (1-\alpha) \mu_D, \quad \bar{\mu}(D) = (1-\beta) \mu_U + \beta \mu_D
\end{equation}

\section{Metodología de Resolución}
Para optimizar el comportamiento del agente, aplicamos la **Ecuación de Bellman de Optimalidad**, permitiendo evaluar decisiones no solo por su retorno inmediato, sino por su valor acumulado futuro:
\begin{equation}
    v^*(s) = \max_{a} \left[ r(s,a) + \gamma \sum_{s' \in S} P(s'|s,a) v^*(s') \right]
\end{equation}

\bibliographystyle{plain}
\begin{thebibliography}{9}
    \bibitem{sutton} Sutton, R. S., \& Barto, A. G. (2018). \textit{Reinforcement learning: An introduction}. MIT press.
\end{thebibliography}

\end{document}
